\lecture{8}{Mon 03 Apr 2023 16:40}{Lebesgue Differentiation Theorem}

\subsection{Lebesgue Differentiation Theorem}

\begin{theorem}[Lebesgue Differentiation Theorem]
	If $f \in L^1[a,b]$, and  $F(x) = \int _a^x f$, then  $F$ is differentiable a.e. and $F' = f$ a.e.
\end{theorem}
\begin{proof}
	From hw, $F$ continuous. WLOG, $f \ge 0$, since we can split into positive and negative parts. Hence $F$ increasing.

	Then by Monotone Differentiation Theorem, $F$ differentiable a.e. and $\int _a^x F' \le F(x) - F(a) = \int _a^x f$ a.e. Hence for all $x$, $\int _a^x (F' - f) \le 0$.

	In fact, $\int _t^x (F' - f) \le 0$ if $t < x$. For, let $G(y) = \int _t^y f = F(y) - F(t)$, replacing $a$ by $t$ above, we get
	\[
		0 \ge \int _t^x \left( G' - f \right) =  \int _t^x (F' - f)
	.\] 
	(Note: need to establish beforehand that differentiation for monotone increasing functions is linear).

	Thus $\int _E (F' - f) \le 0$ for any measurable  $E \subset [a,b]$. Let $E = \{x \in [a,b]: F'(x) \ge f(x)\} $. Then $\int _E (F' - f) = 0$. Hence $F' \le f$ a.e. on $E$. Also $F' < f$ on $[a,b] \setminus E$. Hence $F' \le f$ a.e.

	Now we show the other side of inequality. When $f$ is bounded, let $f \le M \in \mathbb{R}$. Let $g = M - f  \ge 0$. Then
	\[
		G(x) = \int _a^x g = M(x-a) - F(x)
	.\] 
	By what we proved, $G' \le g$ a.e. Hence
	\[
		M - F' \le  M - f \implies f \le  F' 
	.\] 

	Now consider unbounded $f$. Let $f_n(x) = \min\left( f(x), n \right) $. Apply the bounded case to each $f_n$ and apply upward monotone convergence. Specifically, note that $F(x) - F_n(x) = \int_a^x (f - f_n)$ which increases with  $x$. Hence
	\[
		F'(x) \ge F_n'(x) = f_n(x) \to f(x)
	.\] 
	Hence $F'(x) \ge f(x)$ a.e.
\end{proof}


\subsection{Condition for Differentiability}

\begin{definition}[Absolute Continuity]
	$F: [a,b] \to \mathbb{R}$ is absolutely continuous if for all $\varepsilon>0$ there exists $\delta>0$ such that whenever $(x_i, y_i) \subset [a,b]$ are disjoint, $i = 1, \ldots, n$ such that $\sum_{i=1}^{n} (y_i - x_i) < \delta $ implies $\sum_{i=1}^{n} |F(x_i) - F(y_i)| < \varepsilon$.
\end{definition}
\begin{remark}
	\begin{enumerate}
		\item 
	Absolute continuity implies uniform continuity (take $n = 1$ ).
\item Absolute continuity also implies bounded variation.
\item If $g$ increasing, then $m_g$ is a Lebesgue-Steltjes measure (on $\mathcal{B}$, Borel sets). Then $g$ is absolutely continuous iff $m_g$ is absolutely continuous wrt Lebesgue measure ( $m_g \ll m$).
	\end{enumerate}
\end{remark}

\begin{proposition}(Compatibility of definitions of \logseq{Absolute Continuity})
If $g$ increasing, then $m_g$ is a Lebesgue-Steltjes measure (on $\mathcal{B}$, Borel sets). Then $g$ is absolutely continuous iff $m_g$ is absolutely continuous wrt Lebesgue measure ( $m_g \ll m$).
\end{proposition}
\begin{proof}
	Assume $m_g \ll m$. Pick $\varepsilon>0$. Then there exists $\delta>0$ such that $m_g(E) < \varepsilon$ if $m(E) < \delta$, for all Borel  $E \subset [a,b]$. Take $(x_i, y_i)$ as in definition of absolute continuity of $g$. 
	We have 
	\[
		m_g [x,y] = |[x,y]|_g \ge  g(y) - g(x)
	.\] 
	Then
		\begin{align*}
			\sum_{i=1}^{n} |g(y_i) - g(x_i)| 
			&\le m_g[x_i, y_i] \\
			&= m_g(E) < \varepsilon
		.\end{align*}

	The other direction is homework.
\end{proof}

\begin{theorem}[a.e. Differentiability equivalent to Absolute Continuity]
	For $F:[a,b] \to \mathbb{R}$, then $F$ absolutely continuous on $[a,b]$ iff there exists $f$ absolutely integrable such that $F(x) = F(a) + \int _a^x f$.
	
\end{theorem}
\begin{proof}
	If $F'$ absolutely integrable, then by HW, $F'$ is uniformly integrable, i.e. for all $\varepsilon>0$ there exists $\delta>0$ such that $\int _E |F'| < \varepsilon$ when $m(E) < \delta$. Now, for $x < y$, we have
	\[
		\int _x^y F' = \int _a ^y F' - \int _a^x F = F|_x^y
	.\] 
	Hence, if $(x_i, y_i) \subset [a,b]$ as in definition, then for $E = \cup (x_i, y_i)$,
	\[
		\sum |F(y_i) - F(x_i)| = \sum_i |\int _{x_i}^{y_i} F'| \le \sum_i \int_{x_i}^{y_i}|F'| = \int_E |F'| < \varepsilon
	.\] 


	On the other direction, $F$ absolutely continuous implies $F$ is of bounded variation (by HW), which implies $F'$ exists almost everywhere and integrable on $[a,b]$.

	Let $G(x) = \int _a^x F'$. By Lebesgue Differentiation Theorem, $G'(x) = F'(x)$ a.e., i.e. $(G - F)' = 0$ a.e. THe proof of the first part gives $G$ is absolutely continuous. The theorem follows from the next lemma.

	
\end{proof}
\begin{lemma}
		If $g$ is absolutely continuous and $g' = 0$ a.e., then $g$ is constant.
	\end{lemma}
\begin{proof}
	Need to prove $g(a) = g(c)$ for any $c \in (a, b]$. Let $E = \{x \in [a,c]: g'(x) = 0\} $. Then $m([a,c] \setminus E) =0$.

	Let $\mathcal{F}_\varepsilon = \{J \subset [a,c]: J \text{ interval, } J \text{ closed}, [g]_J < \varepsilon |J|\} $, where $[g]_{[p,q]} = |g(q) - g(p)|$. Then  $F_\varepsilon$ is a Vitali cover of $E$: since $g'(x)$ exists at some $x$, for all closed enough $y$ we have $\left| \frac{g(y) - g(x)}{y - x} \right| <\varepsilon$.

	Now apply Vitali's Theorem, there exists disjoint $J_1, J_2, \ldots \in \mathcal{F}_\varepsilon$ such that $m\left( E \setminus  \cup _k J_k \right) = 0$. Since $g$ absolutely continuous, there exists $\delta>0$ for our $\varepsilon>0$ as in the definition. Choose $n$ such that $m\left( E \setminus  \bigcup_{k \le n} J_k \right) < \delta$. Also, $[a,c] \setminus \bigcup_{k \le n} J_k$ is a union of intervals $\bigcup_{k = 1} ^m I_k$.

	Then
	\begin{align*}
		|g(c) - g(a)|
		&\le  \sum_{k=1}^{n} [g]_{J_k} + \sum_{k=1}^{m} [g]_{I_k}\\
		&\le  \sum_{k=1}^{n} \varepsilon |J_k| + \varepsilon\\
		&\le \varepsilon (c-a + 1)
	.\end{align*}
\end{proof}

\subsection{Summary of theory of differentiation}
\begin{enumerate}
	\item If $F$ is of bounded variation, then $F'$ exists a.e., and $F'$ integrable. If $F$ increases, then $\int _a^b F' \le  F|_a^b$.
	\item If $f$ integrable, then  $F(x) = \int _a^x f$ satisfies $F' = f$.
	\item If $F$ is absolutely continuous, then $F(x) = F(a) + \int _a^x F'$.
\end{enumerate}

